% related-work.tex -- Data Availability Committee Paper (RU-V6)

\section{Related Work}

We organize related work into three categories: production DAC systems,
academic data availability protocols, and cryptographic primitives for
secret sharing and signature aggregation.

\subsection{Production DAC Systems}

\subsubsection{StarkEx Validium DAC}

StarkWare's StarkEx~\cite{starkware2024da} pioneered the validium model in
2020, operating in three DA modes: ZK-Rollup (on-chain), Validium (off-chain
with DAC), and Volition (per-vault choice). DAC members run a Committee
Service that receives complete batch data from the operator, validates it
against the STARK proof, and returns ECDSA signatures. The ImmutableX instance
uses a 7-member committee with 5-of-7 threshold; ApeX uses 5 members with
3-of-5~\cite{l2beat2024da}. Each member sees all data---no secret sharing or
erasure coding is employed.

\subsubsection{Polygon CDK DAC}

Polygon CDK~\cite{polygon2024cdk} implements a DAC for its validium mode via
the \texttt{CDKDataCommittee.sol} smart contract~\cite{polygon2024cdknode}.
Committee members store complete batch data. The contract verifies a
configurable M-of-N ECDSA multi-signature with sorted addresses to prevent
duplicates. Members are registered with URLs for data retrieval. As with
StarkEx, each node stores the full dataset.

\subsubsection{Arbitrum AnyTrust (Nova)}

Arbitrum Nova~\cite{arbitrum2024anytrust} introduces the AnyTrust protocol
with a distinctive $(N-1)$-of-$N$ threshold model. The system is secure as
long as at least 2 of $N$ members are honest (honest minority). Committee
members produce BLS-aggregated Data Availability Certificates (DACerts). If
the DAC fails to attest, the system falls back to posting full data on-chain
(rollup mode), providing a strict safety improvement over pure validium. Nova
operates with 6 members and 5-of-6 threshold~\cite{l2beat2024da}. Despite the
superior security model, each member still receives the complete batch data.

\subsubsection{EigenDA}

EigenDA~\cite{eigenlayer2025eigenda, eigenlayer2025v2} is a dedicated DA layer
built on EigenLayer restaking. V2 achieves 100~MB/s write throughput with
5-second average confirmation latency across 14 validators. Data is
Reed-Solomon encoded with 8$\times$ redundancy (any 1,024 of 8,192 chunks
suffice for reconstruction) and committed via KZG polynomial commitments.
Individual operators see only their assigned chunks ($O(1/n)$ of total data),
providing partial privacy. However, chunks contain partial plaintext data, and
the Disperser sees the complete blob before encoding---this is computational
privacy at best, not information-theoretic.

\subsubsection{Celestia}

Celestia~\cite{celestia2024da, celestia2024gosquare} uses Data Availability
Sampling (DAS) rather than a committee. Block data is arranged in a $k \times
k$ matrix, extended to $2k \times 2k$ via 2D Reed-Solomon encoding. Light
nodes randomly sample shares; with $s$ samples, the false acceptance
probability is $<2^{-s}$~\cite{albassam2018fraud}. All data is public---validators
and light nodes see sampled shares in the clear---making it incompatible with
enterprise confidentiality requirements.

\subsection{Privacy Comparison of Production Systems}

Table~\ref{tab:production} summarizes the critical architectural differences
across production DAC systems. The privacy column is the central contribution
of this work: no existing system provides information-theoretic guarantees.

\begin{table*}[t]
\centering
\caption{Production Data Availability Systems: Architecture and Privacy Comparison}
\label{tab:production}
\begin{tabular}{lcccccc}
\toprule
\textbf{System} & \textbf{Members} & \textbf{Threshold} & \textbf{Signature} &
\textbf{Privacy Model} & \textbf{Fallback} & \textbf{L2Beat Rating} \\
\midrule
StarkEx (ImmutableX) & 7 & 5/7 & ECDSA & None (full data) & None & WARNING \\
StarkEx (ApeX)       & 5 & 3/5 & ECDSA & None (full data) & None & BAD \\
Polygon CDK          & Config. & M/N & ECDSA & None (full data) & None & BAD \\
Arbitrum Nova        & 6 & 5/6 & BLS   & None (full data) & Rollup mode & WARNING \\
EigenDA V2           & 200+ & 55\% stake & BLS & Partial (chunks) & None & -- \\
Celestia             & 100+ & 2/3 consensus & Ed25519 & None (public) & N/A & -- \\
Espresso (HotShot)   & All & 1/4 recovery & N/A & Partial (VID) & None & -- \\
\midrule
\textbf{Shamir-DAC (ours)} & \textbf{3} & \textbf{2/3} & \textbf{ECDSA} &
\textbf{Information-theoretic} & \textbf{On-chain DA} & \textbf{--} \\
\bottomrule
\end{tabular}
\end{table*}

\subsection{Academic Data Availability Protocols}

Al-Bassam et al.~\cite{albassam2018fraud} established the theoretical
foundation for DAS, proving that probabilistic sampling can replace
honest-majority assumptions for data availability verification. With $s$
random samples and $>50\%$ data unavailable, false acceptance probability
is $<(1/2)^s$. Hall-Andersen et al.~\cite{hallandersen2024foundations}
provided the first formal cryptographic definitions for DAS as a primitive,
establishing connections between DAS and erasure codes and constructing
schemes without trusted setup.

Nazirkhanova et al.~\cite{nazirkhanova2022semiavidpr} proposed Semi-AVID-PR,
the first verifiable information dispersal protocol specifically designed for
validium rollups that provides privacy against honest-but-curious storage
nodes. Their benchmark achieved $<$3~seconds for 22~MB across 256 nodes with
3.2$\times$ storage expansion. The privacy guarantee is computational (based
on encryption), not information-theoretic. Alhaddad et
al.~\cite{alhaddad2022avid} optimized AVID to near-optimal communication
complexity: $O(|M| + kn^2)$ dispersal with $O(|M|/n + k)$ per-node storage.

Gentry et al.~\cite{gentry2022pvss} demonstrated practical non-interactive
publicly verifiable secret sharing (PVSS) scaling to thousands of parties
based on LWE encryption. While providing public verifiability (anyone can
verify shares were correctly generated), the scheme targets large committees;
for small enterprise DACs ($n \leq 10$), standard Shamir suffices.

\subsection{Cryptographic Primitives}

\subsubsection{Shamir's Secret Sharing}

Shamir~\cite{shamir1979share} introduced the $(k,n)$-threshold scheme:
a secret $S$ is encoded as the constant term of a random polynomial
$f(x) = S + a_1 x + \cdots + a_{k-1} x^{k-1}$ over $\mathrm{GF}(p)$,
and shares $(i, f(i))$ are distributed to $n$ parties. Any $k$ shares
reconstruct $S$ via Lagrange interpolation; fewer than $k$ shares reveal
\emph{zero information} about $S$ (information-theoretic security).
Computational complexity is $O(kn)$ for share generation and $O(k^2)$
for reconstruction.

\subsubsection{BLS Signature Aggregation}

Boneh et al.~\cite{boneh2018compact} introduced accountable-subgroup
multi-signatures (ASM) from BLS, achieving $O(\kappa)$-bit aggregated
signatures independent of signer count. Burdges et
al.~\cite{burdges2022aggregatable} extended this with pairing-free
individual verification. BLS aggregation is used by Arbitrum Nova and
EigenDA. For small committees ($n \leq 10$), ECDSA multi-signature is
simpler and natively supported by EVM via \texttt{ecrecover}, making it
the preferred choice for enterprise DACs on Avalanche Subnet-EVM.

\subsubsection{Poseidon Hash Function}

Grassi et al.~\cite{grassi2021poseidon} designed Poseidon as a hash
function operating natively over prime fields, achieving up to 8$\times$
fewer ZK circuit constraints than Pedersen hashing. Operating over the
BN128 scalar field, Poseidon is the natural choice for data commitment
hashing within the Basis Network validium architecture, ensuring
compatibility with existing ZK circuit infrastructure.
